\cleardoublepage
\chapter*{Abstract}

    Continuous monitoring of cardiac signals, especially the electrocardiogram (ECG), is essential to detect risk situations in patients with heart conditions, particularly elderly individuals engaging in supervised physical activities. This work presents the development of a portable, low-cost system for ECG acquisition and remote monitoring using LoRaWAN technology, an efficient, long-range solution for communication in environments with limited infrastructure. The embedded device was built around the XIAO ESP32-S3 microcontroller, the ADS1293 analog front-end (CJMCU-1293 module) and the Wio-SX1262 LoRa module. To overcome the bandwidth and duty-cycle constraints of LoRaWAN, an edge-computing strategy was adopted: arrhythmia detection is performed on the device itself, and only a compact summary of each ECG segment is transmitted. The on-device detection algorithm was reused from a previously validated firmware provided by its authors; the contribution of this work lies in the signal acquisition, the detector integration, the network infrastructure and the transmission of the events over the LoRaWAN network. The network infrastructure was deployed with a local ChirpStack server running on a Raspberry~Pi~3 with a Radioenge gateway; an ingestion service forwards the events to a cloud database, from which a mobile application displays the detected cardiac events to healthcare professionals in real time. The results validated the network infrastructure, the end-to-end data flow and the edge-detection strategy as a feasible solution for cardiac monitoring over low-power, long-range networks. A sequence of validation batteries with a signal generator characterized the detection performance over both synthetic stimuli and real electrocardiograms from the annotated MIT-BIH database; a paired comparison between the isolated detector and the full acquisition chain showed that acquisition does not degrade the algorithm's decision. Robustness was further assessed under synthetic noise, identifying power-line interference rejection --- a task of the analog front-end --- as the main limitation to be improved.

\textbf{Keywords:} ECG, LoRaWAN, IoT, Remote monitoring, Edge computing, Portable device

\cleardoublepage
\chapter*{Resumo}

O monitoramento contínuo de sinais cardíacos, especialmente do eletrocardiograma (ECG), é essencial para detectar situações de risco em pacientes com problemas cardíacos, particularmente em idosos que realizam atividades físicas supervisionadas. Este trabalho apresenta o desenvolvimento de um sistema portátil e de baixo custo para a aquisição de sinais de ECG e o monitoramento remoto por meio da tecnologia LoRaWAN, uma solução eficiente e de longo alcance para comunicação em ambientes com infraestrutura limitada. O dispositivo embarcado foi construído com base no microcontrolador XIAO ESP32-S3, no \textit{front-end} analógico ADS1293 (módulo CJMCU-1293) e no módulo LoRa Wio-SX1262. Para contornar as restrições de largura de banda e de \textit{duty cycle} do LoRaWAN, adotou-se uma estratégia de processamento na borda: a detecção de arritmias é realizada no próprio dispositivo, e apenas um resumo compacto de cada segmento de ECG é transmitido. O algoritmo de detecção embarcado foi reutilizado de um \textit{firmware} previamente validado, cedido por seus autores (\citeay{moraesFirmwareECG}); a contribuição deste trabalho recai sobre a aquisição do sinal, a integração do detector, a infraestrutura de rede e a transmissão dos eventos via LoRaWAN. A infraestrutura de rede foi implantada com um servidor \textit{ChirpStack} local, executado em um Raspberry~Pi~3 com \textit{gateway} Radioenge; um serviço de ingestão encaminha os eventos a um banco de dados em nuvem, a partir do qual um aplicativo móvel exibe, em tempo real, os eventos cardíacos detectados aos profissionais de saúde. Os resultados validaram a infraestrutura de rede, o fluxo de dados de ponta a ponta e a estratégia de detecção na borda como solução viável para o monitoramento cardíaco sobre redes de baixo consumo e longo alcance. Uma sequência de baterias de validação com gerador de sinais caracterizou quantitativamente o desempenho de detecção, tanto sobre estímulos sintéticos quanto sobre eletrocardiogramas reais da base anotada MIT-BIH; a comparação pareada entre o detector isolado e a cadeia de aquisição completa evidenciou que a aquisição não degrada a decisão do algoritmo. A robustez do sistema foi ainda avaliada sob ruído sintético, identificando-se na rejeição da interferência da rede elétrica — atribuição do estágio analógico de entrada — a principal limitação a aprimorar.

\textbf{Palavras-chave:} ECG, LoRaWAN, IoT, Monitoramento remoto, Computação na borda, Dispositivo portátil
